\documentclass[11pt,a4paper]{article}

\usepackage[utf8]{inputenc}
\usepackage[T1]{fontenc}
\usepackage[french]{babel}
\usepackage{lmodern}
\usepackage{geometry}
\geometry{margin=2cm}
\usepackage{amsmath,amssymb,amsthm}
\usepackage[most]{tcolorbox}
\usepackage{xcolor}
\usepackage{enumitem}
\usepackage{fancyhdr}
\usepackage{array,tabularx}
\usepackage{multicol}
\usepackage{tikz}
\usepackage{hyperref}

\hypersetup{
  colorlinks=true,
  linkcolor=blue!50!black,
  urlcolor=blue!50!black
}

\definecolor{bleu}{RGB}{25,70,140}
\definecolor{vert}{RGB}{30,120,80}
\definecolor{rouge}{RGB}{170,45,45}
\definecolor{orange}{RGB}{210,120,30}
\definecolor{violet}{RGB}{100,70,160}
\definecolor{gris}{RGB}{245,247,250}

\pagestyle{fancy}
\fancyhf{}
\lhead{\textbf{Mathématiques 3e}}
\rhead{\textbf{Nombres, calculs et équations}}
\cfoot{\thepage}

\setlength{\parindent}{0pt}
\setlength{\parskip}{0.45em}

\newcounter{exercice}
\newcounter{correction}

\tcbset{
  enhanced,
  breakable,
  sharp corners,
  boxrule=0.8pt,
  colback=white,
  fonttitle=\bfseries
}

\newtcolorbox{definitionbox}[1][]{
  colframe=bleu,
  colback=blue!3!white,
  title=Définition,
  #1
}

\newtcolorbox{propriete}[1][]{
  colframe=vert,
  colback=green!3!white,
  title=Propriété,
  #1
}

\newtcolorbox{methode}[1][]{
  colframe=violet,
  colback=violet!4!white,
  title=Méthode,
  #1
}

\newtcolorbox{exemple}[1][]{
  colframe=orange,
  colback=orange!4!white,
  title=Exemple corrigé,
  #1
}

\newtcolorbox{attention}[1][]{
  colframe=rouge,
  colback=red!3!white,
  title=Erreur classique / Attention,
  #1
}

\newtcolorbox{exercice}[1][]{
  colframe=black!65,
  colback=gris,
  title={Exercice \refstepcounter{exercice}\theexercice\quad #1},
}

\newtcolorbox{correction}[1][]{
  colframe=vert!70!black,
  colback=green!2!white,
  title={Correction \refstepcounter{correction}\thecorrection\quad #1},
}

\newcommand{\R}{\mathbb{R}}
\newcommand{\Q}{\mathbb{Q}}
\newcommand{\N}{\mathbb{N}}
\newcommand{\Z}{\mathbb{Z}}
\newcommand{\encadre}[1]{\boxed{\displaystyle #1}}

\begin{document}

\begin{titlepage}
\centering
\vspace*{2cm}
{\Huge \textbf{Cours complet de mathématiques}}\\[0.5cm]
{\Large Niveau troisième générale}\\[1.5cm]
{\Huge \textbf{Cours 1}}\\[0.4cm]
{\LARGE \textbf{Nombres, calcul numérique, calcul littéral et équations}}\\[1.5cm]
\begin{tcolorbox}[colframe=bleu,colback=blue!3!white,width=0.9\textwidth]
Ce cours regroupe les bases numériques et algébriques indispensables en troisième : nombres relatifs et rationnels, fractions, puissances, racines carrées, arithmétique, calcul littéral, identités remarquables, équations, inéquations et mise en équation de problèmes.
\end{tcolorbox}
\vfill
{\Large Année scolaire 2026}
\end{titlepage}

\tableofcontents
\newpage

\section*{Préambule : comment travailler ce chapitre ?}
\addcontentsline{toc}{section}{Préambule : comment travailler ce chapitre ?}

Ce chapitre est l'un des plus importants de l'année de troisième. Il sert à calculer, simplifier, raisonner, modéliser et résoudre des problèmes. Il prépare directement le brevet, mais aussi l'entrée en seconde.

\begin{methode}
Pour progresser dans ce chapitre, il faut :
\begin{enumerate}
  \item connaître les définitions et propriétés ;
  \item savoir refaire les exemples corrigés sans regarder ;
  \item s'entraîner sur des calculs courts et réguliers ;
  \item rédiger les problèmes avec des phrases ;
  \item vérifier les résultats obtenus.
\end{enumerate}
\end{methode}

\begin{attention}
Une erreur fréquente consiste à vouloir aller trop vite. En calcul numérique et littéral, les étapes intermédiaires sont essentielles. Une ligne bien écrite vaut souvent mieux qu'un calcul mental risqué.
\end{attention}

\newpage

\section{Nombres relatifs et rationnels}

\subsection{Nombres positifs, négatifs et opposés}

\begin{definitionbox}
Un \textbf{nombre relatif} est un nombre qui peut être positif, négatif ou nul.

\begin{itemize}
  \item Un nombre positif est supérieur ou égal à $0$ : $3$, $7{,}5$, $\frac{2}{3}$.
  \item Un nombre négatif est inférieur ou égal à $0$ : $-4$, $-2{,}8$, $-\frac{5}{7}$.
  \item Le nombre $0$ est à la fois positif et négatif.
\end{itemize}
\end{definitionbox}

\begin{definitionbox}
Deux nombres sont \textbf{opposés} lorsqu'ils ont la même distance à zéro mais des signes contraires.

L'opposé de $a$ est $-a$.

Exemples :
\[
\text{l'opposé de } 7 \text{ est } -7,
\qquad
\text{l'opposé de } -4 \text{ est } 4.
\]
\end{definitionbox}

\begin{exemple}
Dans un relevé de température :
\[
-3^\circ\text{C}
\]
signifie que la température est $3^\circ$C en dessous de zéro.

Si la température passe de $-3^\circ$C à $2^\circ$C, elle augmente de :
\[
2-(-3)=2+3=5^\circ\text{C}.
\]
\end{exemple}

\subsection{Droite graduée, comparaison et distance à zéro}

\begin{definitionbox}
La \textbf{distance à zéro} d'un nombre est sa distance au nombre $0$ sur la droite graduée. Elle est toujours positive.

Exemples :
\[
\text{distance à zéro de } -5 = 5,
\qquad
\text{distance à zéro de } 8 = 8.
\]
\end{definitionbox}

\begin{center}
\begin{tikzpicture}[scale=0.8]
\draw[->] (-6,0)--(6.5,0);
\foreach \x in {-5,-4,-3,-2,-1,0,1,2,3,4,5,6}
{
  \draw (\x,0.1)--(\x,-0.1);
  \node[below] at (\x,-0.1){\x};
}
\draw[very thick,blue] (-4,0.25)--(0,0.25);
\node[above,blue] at (-2,0.25){distance $4$};
\end{tikzpicture}
\end{center}

\begin{propriete}
Sur une droite graduée horizontale, les nombres sont rangés dans l'ordre croissant de gauche à droite.

Ainsi :
\[
-5<-2<0<3<7.
\]
\end{propriete}

\begin{attention}
Pour les nombres négatifs, le nombre qui paraît avoir la plus grande distance à zéro est en réalité le plus petit.

Par exemple :
\[
-12 < -3.
\]
\end{attention}

\subsection{Addition et soustraction de nombres relatifs}

\begin{propriete}
Pour additionner deux nombres relatifs :
\begin{itemize}
  \item s'ils ont le même signe, on garde ce signe et on additionne les distances à zéro ;
  \item s'ils ont des signes contraires, on soustrait les distances à zéro et on garde le signe du nombre qui a la plus grande distance à zéro.
\end{itemize}
\end{propriete}

\begin{exemple}
\[
(-7)+(-4)=-(7+4)=-11.
\]

\[
(-9)+5=-(9-5)=-4.
\]

\[
8+(-3)=8-3=5.
\]
\end{exemple}

\begin{propriete}
Soustraire un nombre revient à ajouter son opposé :
\[
a-b=a+(-b).
\]
\end{propriete}

\begin{exemple}
\[
5-8=5+(-8)=-3.
\]
\[
-4-7=-4+(-7)=-11.
\]
\[
-6-(-10)=-6+10=4.
\]
\end{exemple}

\subsection{Multiplication et division de nombres relatifs}

\begin{propriete}
Règle des signes :
\[
(+)\times(+)=+,
\qquad
(+)\times(-)=-,
\qquad
(-)\times(+)=-,
\qquad
(-)\times(-)=+.
\]

La même règle s'applique à la division.
\end{propriete}

\begin{exemple}
\[
(-3)\times 7=-21.
\]
\[
(-4)\times(-8)=32.
\]
\[
\frac{-18}{3}=-6.
\]
\[
\frac{-24}{-6}=4.
\]
\end{exemple}

\begin{attention}
Le signe $-$ peut avoir deux rôles :
\begin{itemize}
  \item il peut indiquer qu'un nombre est négatif : $-5$ ;
  \item il peut indiquer une soustraction : $8-5$.
\end{itemize}
Il faut donc bien lire l'expression.
\end{attention}

\subsection{Nombres rationnels}

\begin{definitionbox}
Un \textbf{nombre rationnel} est un nombre qui peut s'écrire sous la forme :
\[
\frac{a}{b}
\]
où $a$ et $b$ sont des entiers relatifs et $b\neq 0$.
\end{definitionbox}

\begin{exemple}
Les nombres suivants sont rationnels :
\[
\frac{3}{4},\quad -\frac{7}{5},\quad 2=\frac{2}{1},\quad -6=\frac{-6}{1},\quad 0{,}25=\frac{1}{4}.
\]
\end{exemple}

\begin{attention}
Un nombre rationnel n'a pas toujours une écriture décimale finie. Par exemple :
\[
\frac{1}{3}=0{,}333333\ldots
\]
Son écriture décimale est infinie périodique.
\end{attention}

\begin{exercice}[Application directe]
Calculer :
\begin{multicols}{2}
\begin{enumerate}[label=\alph*)]
  \item $(-5)+(-8)$
  \item $(-12)+7$
  \item $9-14$
  \item $-6-(-11)$
  \item $(-4)\times 9$
  \item $(-7)\times(-3)$
  \item $\dfrac{-30}{5}$
  \item $\dfrac{-42}{-7}$
\end{enumerate}
\end{multicols}
\end{exercice}

\begin{exercice}[Méthode]
Ranger les nombres suivants dans l'ordre croissant :
\[
-3{,}5 \ ;\ 2 \ ;\ -7 \ ;\ 0 \ ;\ -1{,}2 \ ;\ 4{,}1.
\]
\end{exercice}

\begin{exercice}[Problème]
À midi, la température est de $-2^\circ$C. Elle augmente de $7^\circ$C dans l'après-midi puis diminue de $4^\circ$C le soir. Quelle est la température finale ?
\end{exercice}

\newpage

\section{Priorités opératoires}

\subsection{Ordre des opérations}

\begin{propriete}
Dans un calcul, on effectue les opérations dans l'ordre suivant :
\begin{enumerate}
  \item les calculs entre parenthèses ;
  \item les puissances ;
  \item les multiplications et divisions, de gauche à droite ;
  \item les additions et soustractions, de gauche à droite.
\end{enumerate}
\end{propriete}

\begin{exemple}
Calculer :
\[
A=3+4\times 5.
\]
La multiplication est prioritaire :
\[
A=3+20=23.
\]
\end{exemple}

\begin{exemple}
Calculer :
\[
B=(3+4)\times 5.
\]
On commence par les parenthèses :
\[
B=7\times 5=35.
\]
\end{exemple}

\begin{attention}
Les deux calculs précédents donnent des résultats différents :
\[
3+4\times 5=23
\]
mais
\[
(3+4)\times 5=35.
\]
Les parenthèses changent le résultat.
\end{attention}

\subsection{Calculs avec nombres relatifs}

\begin{exemple}
Calculer :
\[
C=-3+5\times(-4).
\]
On commence par la multiplication :
\[
C=-3+(-20).
\]
Donc :
\[
C=-23.
\]
\end{exemple}

\subsection{Calculs avec fractions}

\begin{exemple}
Calculer :
\[
D=\frac{2}{3}+\frac{1}{3}\times 6.
\]
On commence par la multiplication :
\[
\frac{1}{3}\times 6=2.
\]
Donc :
\[
D=\frac{2}{3}+2=\frac{2}{3}+\frac{6}{3}=\frac{8}{3}.
\]
\end{exemple}

\subsection{Utilisation de la calculatrice}

\begin{methode}
Avec la calculatrice, il faut utiliser des parenthèses pour traduire correctement l'expression.

Par exemple, pour calculer :
\[
\frac{3+5}{2}
\]
il faut taper :
\[
(3+5)\div 2.
\]

Si l'on tape seulement :
\[
3+5\div 2,
\]
la calculatrice calcule :
\[
3+\frac{5}{2}.
\]
\end{methode}

\begin{exercice}[Application directe]
Calculer :
\begin{multicols}{2}
\begin{enumerate}[label=\alph*)]
  \item $7+3\times 4$
  \item $(7+3)\times 4$
  \item $20-12\div 3$
  \item $(20-12)\div 4$
  \item $5^2+3$
  \item $2\times 3^2$
  \item $(-4)^2$
  \item $-4^2$
\end{enumerate}
\end{multicols}
\end{exercice}

\begin{exercice}[Méthode]
Calculer en écrivant toutes les étapes :
\[
A=6+2\times(8-5)^2.
\]
\[
B=\left(12-3\times 2\right)\div(-3).
\]
\[
C=\frac{5}{6}+\frac{1}{3}\times 4.
\]
\end{exercice}

\newpage

\section{Fractions}

\subsection{Fraction comme quotient}

\begin{definitionbox}
Une fraction est une écriture de quotient :
\[
\frac{a}{b}
\]
où $a$ est le numérateur, $b$ est le dénominateur et $b\neq 0$.

Elle représente le nombre qui, multiplié par $b$, donne $a$ :
\[
\frac{a}{b}\times b=a.
\]
\end{definitionbox}

\begin{exemple}
\[
\frac{3}{4}
\]
peut représenter trois parts lorsque l'unité a été partagée en quatre parts égales.
\end{exemple}

\subsection{Fractions égales et simplification}

\begin{propriete}
On ne change pas la valeur d'une fraction en multipliant ou en divisant son numérateur et son dénominateur par un même nombre non nul :
\[
\frac{a}{b}=\frac{ka}{kb}
\]
avec $b\neq0$ et $k\neq0$.
\end{propriete}

\begin{exemple}
\[
\frac{12}{18}=\frac{12\div 6}{18\div 6}=\frac{2}{3}.
\]
\end{exemple}

\begin{definitionbox}
Une fraction est \textbf{irréductible} lorsque son numérateur et son dénominateur n'ont pas de diviseur commun autre que $1$.
\end{definitionbox}

\subsection{Addition et soustraction de fractions}

\begin{propriete}
Pour additionner ou soustraire deux fractions, il faut les écrire avec le même dénominateur :
\[
\frac{a}{c}+\frac{b}{c}=\frac{a+b}{c}.
\]
\end{propriete}

\begin{exemple}
\[
\frac{2}{7}+\frac{3}{7}=\frac{5}{7}.
\]
\[
\frac{5}{6}-\frac{1}{6}=\frac{4}{6}=\frac{2}{3}.
\]
\end{exemple}

\begin{exemple}
Calculer :
\[
\frac{3}{4}+\frac{5}{6}.
\]
On cherche un dénominateur commun. On peut choisir $12$ :
\[
\frac{3}{4}=\frac{9}{12},
\qquad
\frac{5}{6}=\frac{10}{12}.
\]
Donc :
\[
\frac{3}{4}+\frac{5}{6}=\frac{9}{12}+\frac{10}{12}=\frac{19}{12}.
\]
\end{exemple}

\subsection{Multiplication de fractions}

\begin{propriete}
Pour multiplier deux fractions, on multiplie les numérateurs entre eux et les dénominateurs entre eux :
\[
\frac{a}{b}\times\frac{c}{d}=\frac{ac}{bd},
\]
avec $b\neq0$ et $d\neq0$.
\end{propriete}

\begin{exemple}
\[
\frac{3}{5}\times \frac{2}{7}=\frac{3\times2}{5\times7}=\frac{6}{35}.
\]
\end{exemple}

\subsection{Division par une fraction}

\begin{definitionbox}
L'inverse d'un nombre non nul $a$ est le nombre qui, multiplié par $a$, donne $1$.

L'inverse de $\frac{a}{b}$ est $\frac{b}{a}$, avec $a\neq0$ et $b\neq0$.
\end{definitionbox}

\begin{propriete}
Diviser par une fraction non nulle revient à multiplier par son inverse :
\[
\frac{a}{b}\div \frac{c}{d}
=
\frac{a}{b}\times \frac{d}{c},
\]
avec $b\neq0$, $c\neq0$ et $d\neq0$.
\end{propriete}

\begin{exemple}
\[
\frac{5}{6}\div\frac{2}{3}
=
\frac{5}{6}\times\frac{3}{2}
=
\frac{15}{12}
=
\frac{5}{4}.
\]
\end{exemple}

\begin{attention}
Pour additionner des fractions, on ne multiplie pas les dénominateurs directement sans réfléchir. Il faut d'abord les mettre au même dénominateur.

Erreur :
\[
\frac{1}{2}+\frac{1}{3}\neq\frac{2}{5}.
\]

En réalité :
\[
\frac{1}{2}+\frac{1}{3}
=
\frac{3}{6}+\frac{2}{6}
=
\frac{5}{6}.
\]
\end{attention}

\begin{exercice}[Application directe]
Simplifier les fractions suivantes :
\begin{multicols}{2}
\begin{enumerate}[label=\alph*)]
  \item $\dfrac{6}{9}$
  \item $\dfrac{15}{25}$
  \item $\dfrac{28}{42}$
  \item $\dfrac{45}{60}$
  \item $\dfrac{36}{48}$
  \item $\dfrac{100}{250}$
\end{enumerate}
\end{multicols}
\end{exercice}

\begin{exercice}[Application directe]
Calculer :
\begin{multicols}{2}
\begin{enumerate}[label=\alph*)]
  \item $\dfrac{2}{5}+\dfrac{1}{5}$
  \item $\dfrac{7}{9}-\dfrac{4}{9}$
  \item $\dfrac{1}{2}+\dfrac{1}{4}$
  \item $\dfrac{3}{5}+\dfrac{2}{3}$
  \item $\dfrac{4}{7}\times\dfrac{3}{5}$
  \item $\dfrac{5}{8}\div\dfrac{3}{4}$
\end{enumerate}
\end{multicols}
\end{exercice}

\begin{exercice}[Problème]
Un gâteau est partagé. Léa mange $\frac{1}{4}$ du gâteau et Sami mange $\frac{1}{3}$ du gâteau.
\begin{enumerate}
  \item Quelle fraction du gâteau ont-ils mangée ensemble ?
  \item Quelle fraction du gâteau reste-t-il ?
\end{enumerate}
\end{exercice}

\newpage

\section{Puissances, exposants négatifs et notation scientifique}

\subsection{Définition d'une puissance}

\begin{definitionbox}
Si $a$ est un nombre et $n$ un entier positif non nul, alors :
\[
a^n=\underbrace{a\times a\times \cdots \times a}_{n\ \text{facteurs}}.
\]

On lit : « $a$ puissance $n$ ».
\end{definitionbox}

\begin{exemple}
\[
2^5=2\times2\times2\times2\times2=32.
\]
\[
(-3)^2=(-3)\times(-3)=9.
\]
\[
-3^2=-(3^2)=-9.
\]
\end{exemple}

\begin{attention}
Il faut distinguer :
\[
(-4)^2=16
\]
et
\[
-4^2=-16.
\]
Dans $(-4)^2$, le nombre mis au carré est $-4$.
Dans $-4^2$, seule la puissance porte sur $4$.
\end{attention}

\subsection{Cas particuliers}

\begin{propriete}
Pour tout nombre $a$ :
\[
a^1=a.
\]

Pour tout nombre non nul $a$ :
\[
a^0=1.
\]
\end{propriete}

\subsection{Puissances de 10 et exposants négatifs}

\begin{propriete}
Pour tout entier positif $n$ :
\[
10^n=1\underbrace{00\ldots0}_{n\ \text{zéros}}.
\]

Et :
\[
10^{-n}=\frac{1}{10^n}.
\]
\end{propriete}

\begin{exemple}
\[
10^3=1000.
\]
\[
10^{-2}=\frac{1}{10^2}=\frac{1}{100}=0{,}01.
\]
\end{exemple}

\subsection{Règles de calcul}

\begin{propriete}
Pour $a\neq0$ et pour tous entiers $m$ et $n$ :
\[
a^m\times a^n=a^{m+n}.
\]
\[
\frac{a^m}{a^n}=a^{m-n}.
\]
\[
(a^m)^n=a^{mn}.
\]
\end{propriete}

\begin{exemple}
\[
2^3\times2^5=2^{3+5}=2^8.
\]
\[
\frac{7^6}{7^2}=7^{6-2}=7^4.
\]
\[
(3^2)^4=3^{2\times4}=3^8.
\]
\end{exemple}

\subsection{Notation scientifique}

\begin{definitionbox}
L'écriture scientifique d'un nombre positif est de la forme :
\[
a\times10^n
\]
où :
\[
1\leq a<10
\]
et $n$ est un entier relatif.
\end{definitionbox}

\begin{exemple}
\[
4500000=4{,}5\times10^6.
\]
\[
0{,}00032=3{,}2\times10^{-4}.
\]
\end{exemple}

\begin{methode}
Pour écrire un nombre en notation scientifique :
\begin{enumerate}
  \item placer la virgule pour obtenir un nombre entre $1$ et $10$ ;
  \item compter le nombre de rangs déplacés ;
  \item utiliser une puissance de $10$ positive si le nombre est grand ;
  \item utiliser une puissance de $10$ négative si le nombre est petit.
\end{enumerate}
\end{methode}

\begin{exercice}[Application directe]
Écrire sous forme de puissance :
\begin{multicols}{2}
\begin{enumerate}[label=\alph*)]
  \item $5\times5\times5$
  \item $(-2)\times(-2)\times(-2)\times(-2)$
  \item $10\times10\times10\times10\times10$
  \item $\dfrac{1}{10^4}$
\end{enumerate}
\end{multicols}
\end{exercice}

\begin{exercice}[Méthode]
Écrire en notation scientifique :
\begin{multicols}{2}
\begin{enumerate}[label=\alph*)]
  \item $350000$
  \item $72000000$
  \item $0{,}0045$
  \item $0{,}00000091$
\end{enumerate}
\end{multicols}
\end{exercice}

\begin{exercice}[Problème]
Une bactérie se divise en deux toutes les 30 minutes. Au départ, il y a une seule bactérie.
\begin{enumerate}
  \item Combien y a-t-il de bactéries au bout d'une heure ?
  \item Combien y a-t-il de bactéries au bout de 5 heures ?
  \item Écrire le résultat sous forme d'une puissance.
\end{enumerate}
\end{exercice}

\newpage

\section{Racines carrées}

\subsection{Définition}

\begin{definitionbox}
Si $a$ est un nombre positif ou nul, la racine carrée de $a$, notée $\sqrt{a}$, est le nombre positif ou nul dont le carré vaut $a$ :
\[
\sqrt{a}=b
\quad \Longleftrightarrow \quad
b^2=a
\quad \text{avec } b\geq0.
\]
\end{definitionbox}

\begin{exemple}
\[
\sqrt{25}=5
\]
car :
\[
5^2=25.
\]
\end{exemple}

\begin{attention}
On écrit :
\[
\sqrt{25}=5
\]
et non pas $-5$, même si :
\[
(-5)^2=25.
\]
La racine carrée désigne toujours le nombre positif ou nul.
\end{attention}

\subsection{Carrés parfaits}

\begin{definitionbox}
Un carré parfait est un nombre qui est le carré d'un entier.
\end{definitionbox}

\[
0^2=0,\quad
1^2=1,\quad
2^2=4,\quad
3^2=9,\quad
4^2=16,\quad
5^2=25,
\]
\[
6^2=36,\quad
7^2=49,\quad
8^2=64,\quad
9^2=81,\quad
10^2=100,\quad
11^2=121,\quad
12^2=144.
\]

\subsection{Valeur exacte et valeur approchée}

\begin{exemple}
Le nombre $\sqrt{17}$ n'est pas un entier. On peut l'encadrer :
\[
4^2=16
\quad \text{et} \quad
5^2=25.
\]
Donc :
\[
4<\sqrt{17}<5.
\]
Avec une calculatrice :
\[
\sqrt{17}\approx4{,}12.
\]
\end{exemple}

\subsection{Simplifier une racine carrée}

\begin{propriete}
Pour $a\geq0$ et $b\geq0$ :
\[
\sqrt{ab}=\sqrt a \times \sqrt b.
\]
Cette propriété permet de simplifier certaines racines.
\end{propriete}

\begin{exemple}
\[
\sqrt{50}=\sqrt{25\times2}=\sqrt{25}\times\sqrt2=5\sqrt2.
\]
\end{exemple}

\subsection{Lien avec la géométrie}

\begin{exemple}
Un carré a une aire de $17\ \text{cm}^2$. Si $c$ est la longueur de son côté, alors :
\[
c^2=17.
\]
Donc :
\[
c=\sqrt{17}\ \text{cm}.
\]
Une valeur approchée est :
\[
c\approx4{,}12\ \text{cm}.
\]
\end{exemple}

\begin{exercice}[Application directe]
Calculer :
\begin{multicols}{2}
\begin{enumerate}[label=\alph*)]
  \item $\sqrt{36}$
  \item $\sqrt{81}$
  \item $\sqrt{100}$
  \item $\sqrt{144}$
  \item $\sqrt{49}$
  \item $\sqrt{121}$
\end{enumerate}
\end{multicols}
\end{exercice}

\begin{exercice}[Méthode]
Simplifier :
\begin{multicols}{2}
\begin{enumerate}[label=\alph*)]
  \item $\sqrt{18}$
  \item $\sqrt{32}$
  \item $\sqrt{45}$
  \item $\sqrt{75}$
\end{enumerate}
\end{multicols}
\end{exercice}

\newpage

\section{Arithmétique}

\subsection{Multiples et diviseurs}

\begin{definitionbox}
Soient $a$ et $b$ deux entiers, avec $a\neq0$.

On dit que $a$ divise $b$ lorsqu'il existe un entier $k$ tel que :
\[
b=a\times k.
\]

On dit aussi que :
\begin{itemize}
  \item $a$ est un diviseur de $b$ ;
  \item $b$ est un multiple de $a$.
\end{itemize}
\end{definitionbox}

\begin{exemple}
$3$ divise $18$ car :
\[
18=3\times6.
\]

Donc $18$ est un multiple de $3$.
\end{exemple}

\subsection{Critères de divisibilité}

\begin{propriete}
Un entier est divisible :
\begin{itemize}
  \item par $2$ s'il est pair ;
  \item par $3$ si la somme de ses chiffres est divisible par $3$ ;
  \item par $5$ s'il se termine par $0$ ou $5$ ;
  \item par $9$ si la somme de ses chiffres est divisible par $9$ ;
  \item par $10$ s'il se termine par $0$.
\end{itemize}
\end{propriete}

\begin{exemple}
Le nombre $378$ est divisible par $3$ car :
\[
3+7+8=18,
\]
et $18$ est divisible par $3$.

Il est aussi divisible par $9$ car $18$ est divisible par $9$.
\end{exemple}

\subsection{Nombres premiers}

\begin{definitionbox}
Un nombre premier est un entier naturel qui possède exactement deux diviseurs positifs : $1$ et lui-même.
\end{definitionbox}

\begin{exemple}
Les premiers nombres premiers sont :
\[
2,\ 3,\ 5,\ 7,\ 11,\ 13,\ 17,\ 19,\ 23,\ 29,\ 31.
\]
\end{exemple}

\begin{attention}
Le nombre $1$ n'est pas premier car il n'a qu'un seul diviseur positif : lui-même.
\end{attention}

\subsection{Décomposition en produit de facteurs premiers}

\begin{propriete}
Tout entier supérieur ou égal à $2$ peut se décomposer en produit de facteurs premiers.
\end{propriete}

\begin{exemple}
Décomposons $60$ :
\[
60=2\times30=2\times2\times15=2^2\times3\times5.
\]
Donc :
\[
60=2^2\times3\times5.
\]
\end{exemple}

\begin{methode}
Pour décomposer un nombre en facteurs premiers :
\begin{enumerate}
  \item on essaie de diviser par $2$ ;
  \item puis par $3$ ;
  \item puis par $5$, $7$, $11$, etc. ;
  \item on continue jusqu'à obtenir uniquement des nombres premiers.
\end{enumerate}
\end{methode}

\subsection{Simplifier une fraction par décomposition}

\begin{exemple}
Simplifier :
\[
\frac{84}{126}.
\]
On décompose :
\[
84=2^2\times3\times7,
\]
\[
126=2\times3^2\times7.
\]
Donc :
\[
\frac{84}{126}
=
\frac{2^2\times3\times7}{2\times3^2\times7}
=
\frac{2}{3}.
\]
\end{exemple}

\subsection{Problèmes de synchronisation}

\begin{exemple}
Deux lampes clignotent. La première s'allume toutes les $12$ secondes, la deuxième toutes les $18$ secondes. Elles s'allument ensemble à midi. Quand s'allumeront-elles de nouveau ensemble ?

On cherche le plus petit multiple commun à $12$ et $18$.

\[
12=2^2\times3,
\qquad
18=2\times3^2.
\]

Le plus petit multiple commun est :
\[
2^2\times3^2=36.
\]

Elles s'allumeront de nouveau ensemble au bout de $36$ secondes.
\end{exemple}

\begin{exercice}[Application directe]
Dire si les affirmations suivantes sont vraies ou fausses :
\begin{enumerate}[label=\alph*)]
  \item $4$ divise $28$.
  \item $7$ divise $45$.
  \item $63$ est un multiple de $9$.
  \item $100$ est divisible par $25$.
\end{enumerate}
\end{exercice}

\begin{exercice}[Méthode]
Décomposer en produit de facteurs premiers :
\begin{multicols}{2}
\begin{enumerate}[label=\alph*)]
  \item $36$
  \item $84$
  \item $120$
  \item $306$
  \item $124$
  \item $2220$
\end{enumerate}
\end{multicols}
\end{exercice}

\begin{exercice}[Problème]
Deux ampoules clignotent. L'une s'allume toutes les $153$ secondes et l'autre toutes les $187$ secondes. À minuit, elles s'allument ensemble.

Au bout de combien de secondes s'allumeront-elles de nouveau ensemble ?
\end{exercice}

\newpage

\section{Calcul littéral}

\subsection{Expression littérale et variable}

\begin{definitionbox}
Une expression littérale est une expression mathématique contenant une ou plusieurs lettres. Ces lettres représentent des nombres.

Exemples :
\[
3x+5,
\qquad
2a-7,
\qquad
x^2+4x+1.
\]
\end{definitionbox}

\begin{definitionbox}
Une variable est une lettre qui peut prendre différentes valeurs.
\end{definitionbox}

\begin{exemple}
Si :
\[
A=3x+5
\]
et si $x=4$, alors :
\[
A=3\times4+5=12+5=17.
\]
\end{exemple}

\subsection{Réduire une expression}

\begin{definitionbox}
Réduire une expression, c'est regrouper les termes de même nature.
\end{definitionbox}

\begin{exemple}
\[
3x+2x=5x.
\]
\[
7a-3a+5=4a+5.
\]
\[
4x^2+3x-2x^2+7x=2x^2+10x.
\]
\end{exemple}

\begin{attention}
On ne peut pas additionner des termes de nature différente.

Par exemple :
\[
3x+2
\]
ne se réduit pas en $5x$.

De même :
\[
x+x=2x
\]
et non pas $x^2$.
\end{attention}

\subsection{Opposé d'une expression}

\begin{propriete}
Prendre l'opposé d'une expression revient à changer tous les signes :
\[
-(a+b)=-a-b.
\]
\[
-(a-b)=-a+b.
\]
\end{propriete}

\begin{exemple}
\[
-(3x-7)=-3x+7.
\]
\[
-(2a+5)=-2a-5.
\]
\end{exemple}

\subsection{Développer}

\begin{definitionbox}
Développer, c'est transformer un produit en somme.
\end{definitionbox}

\begin{propriete}
Distributivité simple :
\[
k(a+b)=ka+kb.
\]
\[
k(a-b)=ka-kb.
\]
\end{propriete}

\begin{exemple}
\[
4(2x-5)=4\times2x-4\times5=8x-20.
\]
\end{exemple}

\begin{propriete}
Double distributivité :
\[
(a+b)(c+d)=ac+ad+bc+bd.
\]
\end{propriete}

\begin{exemple}
Développer :
\[
(2x-3)(5x+7).
\]
On distribue chaque terme :
\[
(2x-3)(5x+7)
=
2x\times5x+2x\times7-3\times5x-3\times7.
\]
Donc :
\[
=10x^2+14x-15x-21.
\]
On réduit :
\[
=10x^2-x-21.
\]
\end{exemple}

\subsection{Factoriser}

\begin{definitionbox}
Factoriser, c'est transformer une somme en produit.
\end{definitionbox}

\begin{methode}
Pour factoriser par facteur commun :
\begin{enumerate}
  \item repérer le facteur présent dans tous les termes ;
  \item le mettre devant une parenthèse ;
  \item écrire ce qu'il reste dans la parenthèse.
\end{enumerate}
\end{methode}

\begin{exemple}
\[
3x+6=3(x+2).
\]
\[
5x^2+10x=5x(x+2).
\]
\end{exemple}

\begin{exercice}[Application directe]
Réduire :
\begin{multicols}{2}
\begin{enumerate}[label=\alph*)]
  \item $3x+5x$
  \item $7a-2a+4$
  \item $5x+3-2x+9$
  \item $4x^2+3x-2x^2+7x$
\end{enumerate}
\end{multicols}
\end{exercice}

\begin{exercice}[Méthode]
Développer et réduire :
\begin{multicols}{2}
\begin{enumerate}[label=\alph*)]
  \item $3(x+5)$
  \item $-2(4x-7)$
  \item $(x+3)(x+5)$
  \item $(2x-1)(3x+4)$
\end{enumerate}
\end{multicols}
\end{exercice}

\begin{exercice}[Méthode]
Factoriser :
\begin{multicols}{2}
\begin{enumerate}[label=\alph*)]
  \item $5x+15$
  \item $7a-21$
  \item $4x^2+8x$
  \item $12x^2-15x$
\end{enumerate}
\end{multicols}
\end{exercice}

\newpage

\section{Identités remarquables}

\subsection{Les trois formules essentielles}

\begin{propriete}
Pour tous nombres $a$ et $b$ :
\[
(a+b)^2=a^2+2ab+b^2.
\]
\[
(a-b)^2=a^2-2ab+b^2.
\]
\[
(a+b)(a-b)=a^2-b^2.
\]
\end{propriete}

\begin{exemple}
\[
(x+5)^2=x^2+2\times x\times5+5^2=x^2+10x+25.
\]
\end{exemple}

\begin{exemple}
\[
(2x-3)^2=(2x)^2-2\times2x\times3+3^2=4x^2-12x+9.
\]
\end{exemple}

\begin{exemple}
\[
(x+7)(x-7)=x^2-49.
\]
\end{exemple}

\subsection{Factoriser avec une identité remarquable}

\begin{methode}
Pour factoriser une différence de deux carrés, on utilise :
\[
a^2-b^2=(a+b)(a-b).
\]
\end{methode}

\begin{exemple}
\[
x^2-49=x^2-7^2=(x+7)(x-7).
\]
\end{exemple}

\begin{exemple}
\[
4x^2-25=(2x)^2-5^2=(2x+5)(2x-5).
\]
\end{exemple}

\begin{attention}
Attention :
\[
a^2+b^2
\]
ne se factorise pas avec l'identité $a^2-b^2$.

Par exemple :
\[
x^2+9
\]
n'est pas égal à $(x+3)(x-3)$, car :
\[
(x+3)(x-3)=x^2-9.
\]
\end{attention}

\begin{exercice}[Application directe]
Développer :
\begin{multicols}{2}
\begin{enumerate}[label=\alph*)]
  \item $(x+4)^2$
  \item $(x-6)^2$
  \item $(x+8)(x-8)$
  \item $(2x+3)^2$
  \item $(3x-5)^2$
  \item $(5x+1)(5x-1)$
\end{enumerate}
\end{multicols}
\end{exercice}

\begin{exercice}[Méthode]
Factoriser :
\begin{multicols}{2}
\begin{enumerate}[label=\alph*)]
  \item $x^2-16$
  \item $x^2-81$
  \item $4x^2-49$
  \item $25x^2-36$
\end{enumerate}
\end{multicols}
\end{exercice}

\newpage

\section{Équations du premier degré}

\subsection{Définition}

\begin{definitionbox}
Une équation est une égalité contenant une inconnue.

Résoudre une équation, c'est trouver toutes les valeurs de l'inconnue qui rendent l'égalité vraie.
\end{definitionbox}

\begin{exemple}
L'équation :
\[
x+3=8
\]
a pour solution :
\[
x=5.
\]
En effet :
\[
5+3=8.
\]
\end{exemple}

\subsection{Principe de la balance}

\begin{propriete}
On peut transformer une équation en effectuant la même opération des deux côtés :
\begin{itemize}
  \item ajouter le même nombre ;
  \item soustraire le même nombre ;
  \item multiplier par le même nombre non nul ;
  \item diviser par le même nombre non nul.
\end{itemize}
\end{propriete}

\begin{exemple}
Résoudre :
\[
3x+5=17.
\]
On soustrait $5$ des deux côtés :
\[
3x=12.
\]
On divise par $3$ :
\[
x=4.
\]
Vérification :
\[
3\times4+5=12+5=17.
\]
Donc la solution est :
\[
\encadre{x=4}.
\]
\end{exemple}

\subsection{Équations avec inconnue des deux côtés}

\begin{exemple}
Résoudre :
\[
4x-8=7x+4.
\]
On soustrait $4x$ des deux côtés :
\[
-8=3x+4.
\]
On soustrait $4$ :
\[
-12=3x.
\]
On divise par $3$ :
\[
x=-4.
\]
Vérification :
\[
4(-4)-8=-16-8=-24,
\]
\[
7(-4)+4=-28+4=-24.
\]
Donc :
\[
\encadre{x=-4}.
\]
\end{exemple}

\subsection{Équations avec parenthèses}

\begin{exemple}
Résoudre :
\[
5(2x-3)=3x+11.
\]
On développe :
\[
10x-15=3x+11.
\]
On soustrait $3x$ :
\[
7x-15=11.
\]
On ajoute $15$ :
\[
7x=26.
\]
On divise par $7$ :
\[
x=\frac{26}{7}.
\]
\end{exemple}

\begin{methode}
Pour résoudre une équation :
\begin{enumerate}
  \item développer si nécessaire ;
  \item réduire chaque membre ;
  \item regrouper les termes en $x$ d'un côté ;
  \item regrouper les nombres de l'autre côté ;
  \item diviser pour isoler $x$ ;
  \item vérifier.
\end{enumerate}
\end{methode}

\begin{exercice}[Application directe]
Résoudre :
\begin{multicols}{2}
\begin{enumerate}[label=\alph*)]
  \item $x+7=15$
  \item $x-9=4$
  \item $3x=21$
  \item $-5x=30$
  \item $2x+3=11$
  \item $7x-4=24$
\end{enumerate}
\end{multicols}
\end{exercice}

\begin{exercice}[Méthode]
Résoudre :
\begin{multicols}{2}
\begin{enumerate}[label=\alph*)]
  \item $4x-8=7x+4$
  \item $5(2x-3)=3x+11$
  \item $3(x+4)=2x-7$
  \item $2(5x-1)-3x=19$
\end{enumerate}
\end{multicols}
\end{exercice}

\newpage

\section{Équations produit}

\subsection{Propriété du produit nul}

\begin{propriete}
Un produit est nul si et seulement si au moins l'un de ses facteurs est nul :
\[
AB=0
\Longleftrightarrow
A=0 \quad \text{ou} \quad B=0.
\]
\end{propriete}

\begin{exemple}
Résoudre :
\[
(2x-7)(3x+4)=0.
\]
D'après la propriété du produit nul :
\[
2x-7=0
\quad \text{ou} \quad
3x+4=0.
\]
Première équation :
\[
2x=7,
\quad
x=\frac{7}{2}.
\]
Deuxième équation :
\[
3x=-4,
\quad
x=-\frac{4}{3}.
\]
Les solutions sont :
\[
\encadre{x=\frac{7}{2}\ \text{ou}\ x=-\frac{4}{3}}.
\]
\end{exemple}

\begin{exemple}
Résoudre :
\[
x(x-5)=0.
\]
Donc :
\[
x=0
\quad \text{ou} \quad
x-5=0.
\]
Ainsi :
\[
x=0
\quad \text{ou} \quad
x=5.
\]
\end{exemple}

\begin{attention}
La propriété du produit nul ne s'applique que lorsque l'équation est sous la forme :
\[
\text{produit}=0.
\]

Par exemple, dans :
\[
(x+3)(x-2)=7,
\]
on ne peut pas écrire directement $x+3=0$ ou $x-2=0$.
\end{attention}

\begin{exercice}[Application directe]
Résoudre :
\begin{multicols}{2}
\begin{enumerate}[label=\alph*)]
  \item $(x-4)(x+2)=0$
  \item $(3x-6)(x+5)=0$
  \item $x(x+9)=0$
  \item $(5x-2)(x+8)=0$
\end{enumerate}
\end{multicols}
\end{exercice}

\newpage

\section{Équations du type \texorpdfstring{$x^2=a$}{x2=a}}

\begin{propriete}
Soit $a$ un nombre réel.

\begin{itemize}
  \item Si $a>0$, alors l'équation $x^2=a$ a deux solutions :
  \[
  x=-\sqrt a
  \quad \text{ou} \quad
  x=\sqrt a.
  \]
  \item Si $a=0$, alors l'équation $x^2=0$ a une seule solution :
  \[
  x=0.
  \]
  \item Si $a<0$, alors l'équation $x^2=a$ n'a pas de solution réelle.
\end{itemize}
\end{propriete}

\begin{exemple}
Résoudre :
\[
x^2=25.
\]
Comme $25>0$, les solutions sont :
\[
x=-\sqrt{25}
\quad \text{ou} \quad
x=\sqrt{25}.
\]
Donc :
\[
\encadre{x=-5\ \text{ou}\ x=5}.
\]
\end{exemple}

\begin{exemple}
Résoudre :
\[
x^2=20.
\]
Comme $20>0$, les solutions exactes sont :
\[
\encadre{x=-\sqrt{20}\ \text{ou}\ x=\sqrt{20}}.
\]
Or :
\[
\sqrt{20}=\sqrt{4\times5}=2\sqrt5.
\]
Donc :
\[
\encadre{x=-2\sqrt5\ \text{ou}\ x=2\sqrt5}.
\]
\end{exemple}

\begin{exemple}
Résoudre :
\[
x^2=-9.
\]
Un carré est toujours positif ou nul. Il n'existe donc aucun nombre réel dont le carré vaut $-9$.

L'équation n'a pas de solution réelle.
\end{exemple}

\begin{exercice}[Application directe]
Résoudre :
\begin{multicols}{2}
\begin{enumerate}[label=\alph*)]
  \item $x^2=36$
  \item $x^2=49$
  \item $x^2=0$
  \item $x^2=-4$
  \item $x^2=12$
  \item $x^2=50$
\end{enumerate}
\end{multicols}
\end{exercice}

\newpage

\section{Inéquations simples}

\subsection{Définition}

\begin{definitionbox}
Une inéquation est une inégalité contenant une inconnue.

Résoudre une inéquation, c'est trouver toutes les valeurs de l'inconnue qui rendent l'inégalité vraie.
\end{definitionbox}

\begin{exemple}
L'inéquation :
\[
x+3<7
\]
est équivalente à :
\[
x<4.
\]
Les solutions sont tous les nombres strictement inférieurs à $4$.
\end{exemple}

\subsection{Règles de résolution}

\begin{propriete}
On peut ajouter ou soustraire le même nombre aux deux membres d'une inégalité sans changer le sens de l'inégalité.
\end{propriete}

\begin{propriete}
On peut multiplier ou diviser les deux membres d'une inégalité par un même nombre strictement positif sans changer le sens de l'inégalité.
\end{propriete}

\begin{propriete}
Lorsqu'on multiplie ou divise les deux membres d'une inégalité par un nombre strictement négatif, on inverse le sens de l'inégalité.
\end{propriete}

\begin{exemple}
Résoudre :
\[
-3x<12.
\]
On divise par $-3$, donc on inverse le sens :
\[
x>-4.
\]
\end{exemple}

\begin{exemple}
Résoudre :
\[
2x-5\geq9.
\]
On ajoute $5$ :
\[
2x\geq14.
\]
On divise par $2$ :
\[
x\geq7.
\]
\end{exemple}

\subsection{Représentation sur une droite graduée}

\begin{center}
\begin{tikzpicture}[scale=0.9]
\draw[->] (-5,0)--(5,0);
\foreach \x in {-4,-3,-2,-1,0,1,2,3,4}
{
  \draw (\x,0.1)--(\x,-0.1);
  \node[below] at (\x,-0.1){\x};
}
\draw[very thick,blue,->] (-4,0.35)--(5,0.35);
\draw[blue,fill=white,thick] (-4,0.35) circle (0.12);
\node[above,blue] at (0,0.45){solutions de $x>-4$};
\end{tikzpicture}
\end{center}

\begin{exercice}[Application directe]
Résoudre :
\begin{multicols}{2}
\begin{enumerate}[label=\alph*)]
  \item $x+5<12$
  \item $x-3\geq 8$
  \item $2x<10$
  \item $-4x\leq 20$
  \item $3x+1>7$
  \item $-2x+5<11$
\end{enumerate}
\end{multicols}
\end{exercice}

\newpage

\section{Problèmes de mise en équation}

\subsection{Méthode générale}

\begin{methode}
Pour résoudre un problème par mise en équation :
\begin{enumerate}
  \item choisir l'inconnue ;
  \item traduire l'énoncé en langage mathématique ;
  \item poser l'équation ;
  \item résoudre l'équation ;
  \item vérifier que la solution est cohérente ;
  \item répondre par une phrase.
\end{enumerate}
\end{methode}

\begin{exemple}
Un nombre augmenté de $7$ vaut $22$. Quel est ce nombre ?

On note $x$ le nombre cherché.

L'énoncé se traduit par :
\[
x+7=22.
\]
On résout :
\[
x=22-7=15.
\]
Le nombre cherché est donc :
\[
\encadre{15}.
\]
\end{exemple}

\subsection{Problème d'âge}

\begin{exemple}
Emma a trois fois l'âge de son petit frère. À eux deux, ils ont $32$ ans. Quel est l'âge de chacun ?

On note $x$ l'âge du petit frère.

Alors l'âge d'Emma est $3x$.

La somme des âges vaut $32$, donc :
\[
x+3x=32.
\]
\[
4x=32.
\]
\[
x=8.
\]

Le petit frère a $8$ ans et Emma a :
\[
3\times8=24
\]
ans.

Réponse : le petit frère a $8$ ans et Emma a $24$ ans.
\end{exemple}

\subsection{Problème de prix}

\begin{exemple}
Un cahier coûte $2$ euros de plus qu'un stylo. Trois cahiers et deux stylos coûtent $19$ euros. Quel est le prix d'un stylo ?

On note $x$ le prix d'un stylo.

Le prix d'un cahier est donc :
\[
x+2.
\]

Le prix total est :
\[
3(x+2)+2x=19.
\]
On développe :
\[
3x+6+2x=19.
\]
On réduit :
\[
5x+6=19.
\]
On soustrait $6$ :
\[
5x=13.
\]
On divise par $5$ :
\[
x=\frac{13}{5}=2{,}60.
\]

Un stylo coûte $2{,}60$ euros.
\end{exemple}

\subsection{Problème de périmètre}

\begin{exemple}
Un rectangle a une longueur qui mesure $5$ cm de plus que sa largeur. Son périmètre vaut $34$ cm. Déterminer ses dimensions.

On note $x$ la largeur.

La longueur vaut :
\[
x+5.
\]

Le périmètre d'un rectangle vaut :
\[
2\times(\text{longueur}+\text{largeur}).
\]

Donc :
\[
2(x+x+5)=34.
\]
\[
2(2x+5)=34.
\]
\[
4x+10=34.
\]
\[
4x=24.
\]
\[
x=6.
\]

La largeur vaut $6$ cm et la longueur vaut :
\[
6+5=11\ \text{cm}.
\]
\end{exemple}

\begin{exercice}[Problème d'âge]
Lucas a deux fois l'âge de sa soeur. Dans $6$ ans, ils auront ensemble $45$ ans. Quel est l'âge actuel de chacun ?
\end{exercice}

\begin{exercice}[Problème de nombres]
La somme de trois nombres entiers consécutifs vaut $84$. Quels sont ces trois nombres ?
\end{exercice}

\begin{exercice}[Problème de prix]
Un cinéma propose deux tarifs :
\begin{itemize}
  \item tarif A : $8$ euros la place ;
  \item tarif B : une carte à $30$ euros puis $5$ euros la place.
\end{itemize}
À partir de combien de places le tarif B devient-il plus avantageux ?
\end{exercice}

\begin{exercice}[Problème géométrique]
Un triangle isocèle a deux côtés égaux mesurant chacun $x$ cm et une base de $8$ cm. Son périmètre vaut $34$ cm. Déterminer $x$.
\end{exercice}

\newpage

\section{Grand entraînement type brevet}

\begin{exercice}[Calcul numérique]
Calculer en détaillant :
\[
A=7-3\times(5-9)^2.
\]
\[
B=\frac{3}{4}+\frac{5}{6}\times\frac{2}{3}.
\]
\[
C=2^5\times2^{-3}.
\]
\end{exercice}

\begin{exercice}[Arithmétique]
\begin{enumerate}
  \item Décomposer $180$ et $252$ en produit de facteurs premiers.
  \item Simplifier la fraction $\dfrac{180}{252}$.
  \item Deux événements se produisent respectivement toutes les $180$ secondes et toutes les $252$ secondes. Ils se produisent ensemble à $10$h. Au bout de combien de temps se reproduiront-ils ensemble ?
\end{enumerate}
\end{exercice}

\begin{exercice}[Calcul littéral]
On considère :
\[
E=(2x-3)^2-(x+5)(x-5).
\]
\begin{enumerate}
  \item Développer et réduire $E$.
  \item Calculer $E$ pour $x=2$.
\end{enumerate}
\end{exercice}

\begin{exercice}[Équations]
Résoudre :
\[
3(2x-5)=4x+7.
\]
\[
(2x-1)(x+6)=0.
\]
\[
x^2=45.
\]
\end{exercice}

\begin{exercice}[Mise en équation]
Un abonnement de sport coûte $40$ euros d'inscription puis $15$ euros par mois. Une autre salle propose $25$ euros par mois sans inscription.

À partir de combien de mois le premier abonnement devient-il moins cher que le second ?
\end{exercice}

\newpage

\section{Corrections détaillées}

\begin{correction}[Nombres relatifs]
\begin{enumerate}[label=\alph*)]
  \item $(-5)+(-8)=-(5+8)=-13$.
  \item $(-12)+7=-(12-7)=-5$.
  \item $9-14=9+(-14)=-5$.
  \item $-6-(-11)=-6+11=5$.
  \item $(-4)\times9=-36$.
  \item $(-7)\times(-3)=21$.
  \item $\dfrac{-30}{5}=-6$.
  \item $\dfrac{-42}{-7}=6$.
\end{enumerate}
\end{correction}

\begin{correction}[Rangement]
Sur une droite graduée, les nombres les plus petits sont à gauche. Les négatifs sont inférieurs aux positifs. On obtient :
\[
\encadre{-7<-3{,}5<-1{,}2<0<2<4{,}1}.
\]
\end{correction}

\begin{correction}[Température]
La température initiale est $-2^\circ$C. Après une hausse de $7^\circ$C :
\[
-2+7=5.
\]
Après une baisse de $4^\circ$C :
\[
5-4=1.
\]
La température finale est donc :
\[
\encadre{1^\circ\text{C}}.
\]
\end{correction}

\begin{correction}[Priorités opératoires]
\begin{enumerate}[label=\alph*)]
  \item $7+3\times4=7+12=19$.
  \item $(7+3)\times4=10\times4=40$.
  \item $20-12\div3=20-4=16$.
  \item $(20-12)\div4=8\div4=2$.
  \item $5^2+3=25+3=28$.
  \item $2\times3^2=2\times9=18$.
  \item $(-4)^2=16$.
  \item $-4^2=-(4^2)=-16$.
\end{enumerate}
\end{correction}

\begin{correction}[Fractions : simplification]
\[
\frac{6}{9}=\frac{2}{3},
\qquad
\frac{15}{25}=\frac{3}{5},
\qquad
\frac{28}{42}=\frac{2}{3}.
\]
\[
\frac{45}{60}=\frac{3}{4},
\qquad
\frac{36}{48}=\frac{3}{4},
\qquad
\frac{100}{250}=\frac{2}{5}.
\]
\end{correction}

\begin{correction}[Fractions : opérations]
\[
\frac{2}{5}+\frac{1}{5}=\frac{3}{5}.
\]
\[
\frac{7}{9}-\frac{4}{9}=\frac{3}{9}=\frac{1}{3}.
\]
\[
\frac{1}{2}+\frac{1}{4}=\frac{2}{4}+\frac{1}{4}=\frac{3}{4}.
\]
\[
\frac{3}{5}+\frac{2}{3}=\frac{9}{15}+\frac{10}{15}=\frac{19}{15}.
\]
\[
\frac{4}{7}\times\frac{3}{5}=\frac{12}{35}.
\]
\[
\frac{5}{8}\div\frac{3}{4}
=
\frac{5}{8}\times\frac{4}{3}
=
\frac{20}{24}
=
\frac{5}{6}.
\]
\end{correction}

\begin{correction}[Gâteau]
Léa et Sami mangent :
\[
\frac{1}{4}+\frac{1}{3}
=
\frac{3}{12}+\frac{4}{12}
=
\frac{7}{12}.
\]
Il reste :
\[
1-\frac{7}{12}
=
\frac{12}{12}-\frac{7}{12}
=
\frac{5}{12}.
\]
Ils ont mangé $\frac{7}{12}$ du gâteau et il reste $\frac{5}{12}$.
\end{correction}

\begin{correction}[Notation scientifique]
\[
350000=3{,}5\times10^5.
\]
\[
72000000=7{,}2\times10^7.
\]
\[
0{,}0045=4{,}5\times10^{-3}.
\]
\[
0{,}00000091=9{,}1\times10^{-7}.
\]
\end{correction}

\begin{correction}[Bactéries]
Toutes les $30$ minutes, le nombre est multiplié par $2$.

En $1$ heure, il y a deux périodes de $30$ minutes :
\[
2^2=4.
\]

En $5$ heures, il y a :
\[
5\times2=10
\]
périodes de $30$ minutes.

Donc il y a :
\[
2^{10}=1024
\]
bactéries.
\end{correction}

\begin{correction}[Racines carrées]
\[
\sqrt{36}=6,
\quad
\sqrt{81}=9,
\quad
\sqrt{100}=10,
\quad
\sqrt{144}=12,
\quad
\sqrt{49}=7,
\quad
\sqrt{121}=11.
\]
\end{correction}

\begin{correction}[Simplification de racines]
\[
\sqrt{18}=\sqrt{9\times2}=3\sqrt2.
\]
\[
\sqrt{32}=\sqrt{16\times2}=4\sqrt2.
\]
\[
\sqrt{45}=\sqrt{9\times5}=3\sqrt5.
\]
\[
\sqrt{75}=\sqrt{25\times3}=5\sqrt3.
\]
\end{correction}

\begin{correction}[Arithmétique]
\begin{enumerate}[label=\alph*)]
  \item $4$ divise $28$ car $28=4\times7$. Vrai.
  \item $7$ ne divise pas $45$ car $45$ n'est pas dans la table de $7$. Faux.
  \item $63$ est un multiple de $9$ car $63=9\times7$. Vrai.
  \item $100$ est divisible par $25$ car $100=25\times4$. Vrai.
\end{enumerate}
\end{correction}

\begin{correction}[Décomposition]
\[
36=2^2\times3^2.
\]
\[
84=2^2\times3\times7.
\]
\[
120=2^3\times3\times5.
\]
\[
306=2\times3^2\times17.
\]
\[
124=2^2\times31.
\]
\[
2220=2^2\times3\times5\times37.
\]
\end{correction}

\begin{correction}[Ampoules]
On décompose :
\[
153=3^2\times17,
\qquad
187=11\times17.
\]
Le plus petit multiple commun est :
\[
3^2\times 11\times17=1683.
\]
Les deux ampoules s'allumeront de nouveau ensemble au bout de :
\[
\encadre{1683\ \text{secondes}}.
\]
\end{correction}

\begin{correction}[Calcul littéral : réduire]
\[
3x+5x=8x.
\]
\[
7a-2a+4=5a+4.
\]
\[
5x+3-2x+9=3x+12.
\]
\[
4x^2+3x-2x^2+7x=2x^2+10x.
\]
\end{correction}

\begin{correction}[Développer]
\[
3(x+5)=3x+15.
\]
\[
-2(4x-7)=-8x+14.
\]
\[
(x+3)(x+5)=x^2+5x+3x+15=x^2+8x+15.
\]
\[
(2x-1)(3x+4)=6x^2+8x-3x-4=6x^2+5x-4.
\]
\end{correction}

\begin{correction}[Factoriser]
\[
5x+15=5(x+3).
\]
\[
7a-21=7(a-3).
\]
\[
4x^2+8x=4x(x+2).
\]
\[
12x^2-15x=3x(4x-5).
\]
\end{correction}

\begin{correction}[Identités remarquables]
\[
(x+4)^2=x^2+8x+16.
\]
\[
(x-6)^2=x^2-12x+36.
\]
\[
(x+8)(x-8)=x^2-64.
\]
\[
(2x+3)^2=4x^2+12x+9.
\]
\[
(3x-5)^2=9x^2-30x+25.
\]
\[
(5x+1)(5x-1)=25x^2-1.
\]
\end{correction}

\begin{correction}[Factorisation par identité remarquable]
\[
x^2-16=(x+4)(x-4).
\]
\[
x^2-81=(x+9)(x-9).
\]
\[
4x^2-49=(2x+7)(2x-7).
\]
\[
25x^2-36=(5x+6)(5x-6).
\]
\end{correction}

\begin{correction}[Équations simples]
\[
x+7=15 \Longleftrightarrow x=8.
\]
\[
x-9=4 \Longleftrightarrow x=13.
\]
\[
3x=21 \Longleftrightarrow x=7.
\]
\[
-5x=30 \Longleftrightarrow x=-6.
\]
\[
2x+3=11 \Longleftrightarrow 2x=8 \Longleftrightarrow x=4.
\]
\[
7x-4=24 \Longleftrightarrow 7x=28 \Longleftrightarrow x=4.
\]
\end{correction}

\begin{correction}[Équations avec transformations]
\[
4x-8=7x+4
\]
\[
-8=3x+4
\]
\[
-12=3x
\]
\[
x=-4.
\]

\[
5(2x-3)=3x+11
\]
\[
10x-15=3x+11
\]
\[
7x=26
\]
\[
x=\frac{26}{7}.
\]
\end{correction}

\begin{correction}[Équations produit]
\[
(x-4)(x+2)=0
\]
donc :
\[
x=4 \quad \text{ou} \quad x=-2.
\]

\[
(3x-6)(x+5)=0
\]
donc :
\[
3x-6=0 \quad \text{ou} \quad x+5=0.
\]
Ainsi :
\[
x=2 \quad \text{ou} \quad x=-5.
\]
\end{correction}

\begin{correction}[Équations $x^2=a$]
\[
x^2=36 \Longleftrightarrow x=-6 \text{ ou } x=6.
\]
\[
x^2=49 \Longleftrightarrow x=-7 \text{ ou } x=7.
\]
\[
x^2=0 \Longleftrightarrow x=0.
\]
\[
x^2=-4
\]
n'a pas de solution réelle.
\[
x^2=12 \Longleftrightarrow x=-\sqrt{12} \text{ ou } x=\sqrt{12}.
\]
Or $\sqrt{12}=2\sqrt3$, donc :
\[
x=-2\sqrt3 \text{ ou } x=2\sqrt3.
\]
\end{correction}

\begin{correction}[Inéquations]
\[
x+5<12 \Longleftrightarrow x<7.
\]
\[
x-3\geq8 \Longleftrightarrow x\geq11.
\]
\[
2x<10 \Longleftrightarrow x<5.
\]
\[
-4x\leq20.
\]
On divise par $-4$, donc on inverse le sens :
\[
x\geq -5.
\]
\[
3x+1>7 \Longleftrightarrow 3x>6 \Longleftrightarrow x>2.
\]
\[
-2x+5<11 \Longleftrightarrow -2x<6.
\]
On divise par $-2$, donc on inverse :
\[
x>-3.
\]
\end{correction}

\begin{correction}[Problème d'âge]
On note $x$ l'âge de la soeur. Lucas a donc $2x$ ans.

Dans $6$ ans, la soeur aura $x+6$ ans et Lucas aura $2x+6$ ans.

Ils auront ensemble $45$ ans :
\[
(x+6)+(2x+6)=45.
\]
\[
3x+12=45.
\]
\[
3x=33.
\]
\[
x=11.
\]
La soeur a $11$ ans et Lucas a :
\[
2\times11=22
\]
ans.
\end{correction}

\begin{correction}[Nombres consécutifs]
On note $x$ le premier entier. Les trois entiers consécutifs sont :
\[
x,\quad x+1,\quad x+2.
\]
Leur somme vaut $84$ :
\[
x+(x+1)+(x+2)=84.
\]
\[
3x+3=84.
\]
\[
3x=81.
\]
\[
x=27.
\]
Les trois nombres sont :
\[
\encadre{27,\ 28,\ 29}.
\]
\end{correction}

\begin{correction}[Tarifs cinéma]
Soit $x$ le nombre de places.

Tarif A :
\[
8x.
\]
Tarif B :
\[
30+5x.
\]

On cherche quand le tarif B est plus avantageux :
\[
30+5x<8x.
\]
\[
30<3x.
\]
\[
10<x.
\]
Donc le tarif B devient plus avantageux à partir de :
\[
\encadre{11\ \text{places}}.
\]
\end{correction}

\begin{correction}[Triangle isocèle]
Le périmètre vaut :
\[
x+x+8=34.
\]
Donc :
\[
2x+8=34.
\]
\[
2x=26.
\]
\[
x=13.
\]
Chaque côté égal mesure :
\[
\encadre{13\ \text{cm}}.
\]
\end{correction}

\begin{correction}[Grand entraînement : calcul numérique]
\[
A=7-3\times(5-9)^2.
\]
\[
A=7-3\times(-4)^2.
\]
\[
A=7-3\times16.
\]
\[
A=7-48=-41.
\]

\[
B=\frac{3}{4}+\frac{5}{6}\times\frac{2}{3}.
\]
\[
B=\frac{3}{4}+\frac{10}{18}.
\]
\[
B=\frac{3}{4}+\frac{5}{9}.
\]
\[
B=\frac{27}{36}+\frac{20}{36}=\frac{47}{36}.
\]

\[
C=2^5\times2^{-3}=2^{5-3}=2^2=4.
\]
\end{correction}

\begin{correction}[Grand entraînement : arithmétique]
\[
180=18\times10=2\times3^2\times2\times5=2^2\times3^2\times5.
\]
\[
252=4\times63=2^2\times3^2\times7.
\]

Donc :
\[
\frac{180}{252}
=
\frac{2^2\times3^2\times5}{2^2\times3^2\times7}
=
\frac{5}{7}.
\]

Le plus petit multiple commun est :
\[
2^2\times3^2\times5\times7=1260.
\]
Les événements se reproduiront ensemble au bout de :
\[
\encadre{1260\ \text{secondes}}.
\]
\end{correction}

\begin{correction}[Grand entraînement : calcul littéral]
\[
E=(2x-3)^2-(x+5)(x-5).
\]
On développe :
\[
(2x-3)^2=4x^2-12x+9.
\]
\[
(x+5)(x-5)=x^2-25.
\]
Donc :
\[
E=4x^2-12x+9-(x^2-25).
\]
\[
E=4x^2-12x+9-x^2+25.
\]
\[
E=3x^2-12x+34.
\]

Pour $x=2$ :
\[
E=3\times2^2-12\times2+34.
\]
\[
E=12-24+34=22.
\]
\end{correction}

\begin{correction}[Grand entraînement : équations]
\[
3(2x-5)=4x+7.
\]
\[
6x-15=4x+7.
\]
\[
2x=22.
\]
\[
x=11.
\]

\[
(2x-1)(x+6)=0.
\]
Donc :
\[
2x-1=0 \quad \text{ou} \quad x+6=0.
\]
Ainsi :
\[
x=\frac12 \quad \text{ou} \quad x=-6.
\]

\[
x^2=45.
\]
Donc :
\[
x=-\sqrt{45} \quad \text{ou} \quad x=\sqrt{45}.
\]
Or :
\[
\sqrt{45}=3\sqrt5.
\]
Donc :
\[
x=-3\sqrt5 \quad \text{ou} \quad x=3\sqrt5.
\]
\end{correction}

\begin{correction}[Grand entraînement : abonnement]
Soit $x$ le nombre de mois.

Premier abonnement :
\[
40+15x.
\]

Deuxième abonnement :
\[
25x.
\]

On cherche :
\[
40+15x<25x.
\]
\[
40<10x.
\]
\[
4<x.
\]
Le premier abonnement devient moins cher à partir de :
\[
\encadre{5\ \text{mois}}.
\]
\end{correction}

\newpage

\section{Fiche de synthèse finale}

\begin{tcolorbox}[colframe=bleu,colback=blue!3!white,title=À retenir absolument]
\begin{multicols}{2}
\begin{itemize}
  \item Soustraire revient à ajouter l'opposé :
  \[
  a-b=a+(-b).
  \]
  \item Règle des signes :
  \[
  (-)\times(-)=+.
  \]
  \item Addition de fractions :
  \[
  \frac{a}{c}+\frac{b}{c}=\frac{a+b}{c}.
  \]
  \item Multiplication :
  \[
  \frac{a}{b}\times\frac{c}{d}=\frac{ac}{bd}.
  \]
  \item Division :
  \[
  \frac{a}{b}\div\frac{c}{d}
  =
  \frac{a}{b}\times\frac{d}{c}.
  \]
  \item Puissances :
  \[
  a^m a^n=a^{m+n}.
  \]
  \item Racine carrée :
  \[
  \sqrt a=b \Longleftrightarrow b^2=a,\ b\geq0.
  \]
  \item Identités remarquables :
  \[
  (a+b)^2=a^2+2ab+b^2.
  \]
  \[
  (a-b)^2=a^2-2ab+b^2.
  \]
  \[
  (a+b)(a-b)=a^2-b^2.
  \]
  \item Produit nul :
  \[
  AB=0\Longleftrightarrow A=0\ \text{ou}\ B=0.
  \]
  \item Si on divise une inégalité par un nombre négatif, on inverse le sens.
\end{itemize}
\end{multicols}
\end{tcolorbox}

\begin{tcolorbox}[colframe=rouge,colback=red!3!white,title=Erreurs classiques à éviter]
\begin{itemize}
  \item Oublier les priorités opératoires.
  \item Confondre $(-4)^2$ et $-4^2$.
  \item Additionner les dénominateurs :
  \[
  \frac12+\frac13\neq\frac25.
  \]
  \item Écrire $x+x=x^2$ au lieu de $x+x=2x$.
  \item Oublier de distribuer le signe moins.
  \item Appliquer la propriété du produit nul quand l'équation n'est pas égale à zéro.
  \item Oublier les deux solutions de $x^2=a$ lorsque $a>0$.
  \item Ne pas inverser le sens d'une inégalité après division par un nombre négatif.
\end{itemize}
\end{tcolorbox}

\end{document}